% Part: incompleteness
% Chapter: representability-in-q
% Section: c

\documentclass[../../../include/open-logic-section]{subfiles}

\begin{document}

\olfileid[hi]{inc}{req}{cee}
\olsection{फलनों का समुच्चय $C$}

$C$ को वह सबसे छोटा फलन-समुच्चय मानें जिसमें
\begin{enumerate}
\item शून्य फलन~$0$,
\item उत्तराधिकारी फलन,
\item जोड़,
\item गुणन,
\item प्रक्षेप, और
\item समता का अभिलाक्षणिक फलन $\Char{=}$
\end{enumerate}
हों और जो निम्न संक्रियाओं के अंतर्गत संवृत हो:
\begin{enumerate}
\item संयोजन, और
\item नियमित फलनों पर लागू अपरिबद्ध खोज।
\end{enumerate}
याद रखें कि अंतिम प्रतिबंध का अर्थ केवल यह है कि आप $\mu$ संक्रिया का
उपयोग तभी कर सकते हैं जब परिणाम सर्वत्र-परिभाषित हो। इसकी तुलना
\emph{सामान्य पुनरावर्ती} फलनों की परिभाषा से करें: यहाँ हमने जोड़, गुणन और
$\Char{=}$ जोड़े हैं, किंतु आदिम पुनरावर्तन हटा दिया है।

स्पष्टतः $C$ का प्रत्येक फलन पुनरावर्ती है, क्योंकि जोड़, गुणन और
$\Char{=}$ पुनरावर्ती हैं। हम दिखाएँगे कि इसका विलोम भी सत्य है; इसका अर्थ
है कि $C$ की शेष सामग्री से हम आदिम पुनरावर्तन कर सकते हैं।

\end{document}
